% =====================  SECTION IV  =====================
\section{Experimental Setup and Results}\label{sec:results}

This section reports what the system was measured to do. To keep the evidential status of
each figure explicit, results are labeled as one of four kinds: an experimental result
(measured on held-out data), a system output (a value the pipeline emits by
construction), a simulation result (an outcome of the deterministic projection), or a
theoretical property (a guarantee that follows from the design). Every quantitative claim
is drawn from the verified metric artifacts.

\subsection{Experimental Setup}\label{subsec:setup}
All models were trained with scikit-learn 1.5.0 under a fixed random seed for
reproducibility, using the configuration in Table~\ref{tab:config}. Headline metrics are
computed on a stratified 20\% holdout of 600 records; stability is estimated with
stratified five-fold cross-validation in which the clustering is refitted inside each
fold so that no proxy-label information crosses the split. Uncertainty on the headline
metrics is quantified by a bootstrap 95\% confidence interval. The proxy-label
distribution is imbalanced---Low 44.9\%, Medium 45.9\%, High 9.2\%---which is the
backdrop against which the minority High-risk results should be read.

\subsection{Clustering Quality}\label{subsec:clustering-quality}
\textbf{\textit{Experimental result.}} The chosen three-cluster solution attains a
silhouette coefficient of 0.370, a Calinski--Harabasz index of 2158.37, and a
Davies--Bouldin index of 0.925 on the standardized ratios, with assignments stable under
resampling at an adjusted Rand index of $0.868 \pm 0.064$. The three clusters map onto
financially distinct, monotonically ordered groups, which supports their use as risk
proxies: as Table~\ref{tab:cluster} shows, mean surplus falls and mean expense ratio
rises steadily from the Low to the High band. Fig.~\ref{fig:cluster} shows the separation
in principal-component space and in surplus--expense-ratio space with the guardrail
thresholds overlaid.

\begin{table}[t]
\centering
\caption{Mean financial profile of each risk cluster (system output)}
\label{tab:cluster}
\begin{tabularx}{\linewidth}{@{}l C C C@{}}
\toprule
Risk band & Mean surplus & Mean expense ratio & Mean debt pressure \\
\midrule
Low & $+1171.67$ & 76.04\% & 10.43\% \\
Medium & $-583.14$ & 117.76\% & 14.30\% \\
High & $-1994.60$ & 195.80\% & 27.54\% \\
\bottomrule
\end{tabularx}
\par\vspace{3pt}
{\footnotesize\itshape\raggedright The debt-pressure column is informative for profiling
but is set to zero on the guardrail path at inference; the clusters themselves are formed
on all three ratios.\par}
\end{table}

\begin{figure}[t]
\centering
\includegraphics[width=\linewidth]{fig_cluster_sep.png}
\caption{K-means separation of the three risk segments (silhouette 0.370). Left: the
first two principal components. Right: the same users in surplus--expense-ratio space with
guardrail thresholds overlaid.}
\label{fig:cluster}
\end{figure}

\subsection{Classification Performance}\label{subsec:classification-perf}
\textbf{\textit{Experimental result.}} On the 600-record holdout the MLP attains 91.0\%
accuracy and a 0.881 macro-F1, with a macro one-vs-rest ROC-AUC of 0.986 and a macro
PR-AUC of 0.951. Five-fold cross-validation gives 91.2\% $\pm$ 0.9\% accuracy and
$0.867 \pm 0.018$ macro-F1, so the holdout figure is representative rather than a
favorable split, and a bootstrap 95\% confidence interval places holdout accuracy in
[0.888, 0.933] and macro-F1 in [0.847, 0.912]. Performance is uneven across bands
(Table~\ref{tab:perclass}): the model is strong on the well-represented Low and Medium
bands but markedly weaker on the minority High band, whose F1 of 0.791 reflects its small
support (62 cases) and its overlap with the Medium band.

\begin{table}[t]
\centering
\caption{Per-class classifier performance of the deployed MLP on the holdout
(experimental result)}
\label{tab:perclass}
\begin{tabularx}{\linewidth}{@{}l C C C C@{}}
\toprule
Risk band & Precision & Recall & F1 score & Support \\
\midrule
Low & 0.922 & 0.971 & 0.946 & 242 \\
Medium & 0.935 & 0.878 & 0.906 & 296 \\
High & 0.761 & 0.823 & 0.791 & 62 \\
\midrule
Macro average & 0.873 & 0.891 & 0.881 & 600 \\
Weighted average & 0.912 & 0.910 & 0.910 & 600 \\
\bottomrule
\end{tabularx}
\par\vspace{3pt}
{\footnotesize\itshape\raggedright Precision, recall, and F1 are computed directly from
the MLP confusion matrix in Fig.~\ref{fig:cm}(a); the macro average is unweighted across
classes and the weighted average is weighted by support. Weighted recall equals overall
accuracy (0.910) by construction.\par}
\end{table}

Fig.~\ref{fig:cm} shows the confusion matrix of every evaluated model on the identical
600-case holdout, which exposes the error structure that the aggregate scores hide. Three
properties hold across all four models. First, errors are overwhelmingly single-band
confusions between adjacent classes. Second---and this matters for a conservative
advisory system---not one model ever confuses Low risk with High risk in either
direction: both off-diagonal corners are exactly zero for the MLP, logistic regression,
random forest, and decision tree alike, so on the classifier path the system never
commits the catastrophic inversion of calling a high-risk user low-risk or vice versa.
Third, the residual errors concentrate in the Medium band and at the Medium--High
boundary, which is precisely where the deterministic guardrail is positioned to intervene.

The minority High band is where the models differ most, and the comparison does not favor
the MLP. Reading recall along the High row of Fig.~\ref{fig:cm}, logistic regression
recovers 61 of 62 true High-risk users (recall 0.984), against 51 of 62 for both the MLP
and the random forest (0.823) and 54 of 62 for the decision tree (0.871). Logistic
regression attains this by escalating more readily---it also relabels 36 Medium users as
High---so its High-band precision falls to 0.629 against the MLP's 0.761, the familiar
precision--recall trade-off on a small class. For a safety-first system that would rather
over- than under-state risk, the higher High-risk sensitivity of logistic regression is
arguably the more desirable operating point, which further weakens any case for preferring
the neural model on the minority class. The per-class ROC and precision--recall curves for
the MLP in Fig.~\ref{fig:rocpr} corroborate the aggregate picture: near-perfect
separability on the well-supported Low band and a visibly lower precision--recall area on
the 62-case High band.

\begin{figure}[t]
\centering
\includegraphics[width=\linewidth]{fig_confusion_all.png}
\caption{Confusion matrices of all four evaluated classifiers on the identical 600-case
holdout (rows: true label; columns: predicted label; each cell shows the count and the
row percentage). Errors are confined to adjacent bands and no model produces any
Low$\leftrightarrow$High confusion. All four panels are computed from the same held-out
predictions and ground-truth proxy labels used for
Tables~\ref{tab:perclass}--\ref{tab:stats}, and reproduce the accuracy and macro-F1 of
Table~\ref{tab:modelcompare} exactly.}
\label{fig:cm}
\end{figure}

\begin{figure}[t]
\centering
\includegraphics[width=\linewidth]{fig_roc_pr.png}
\caption{Per-class one-vs-rest ROC curves (left) and precision--recall curves (right) for
the MLP. The High-risk band shows the lowest average precision.}
\label{fig:rocpr}
\end{figure}

\subsection{Comparative Analysis}\label{subsec:comparative-analysis}
\textbf{\textit{Experimental result.}} The MLP was compared against logistic regression, a
decision tree, and a random forest under the same protocol
(Table~\ref{tab:modelcompare}, Fig.~\ref{fig:modelcompare}). The MLP leads on accuracy and
macro-F1, but the margin over the two strongest baselines is slight, and logistic
regression in fact marginally exceeds the MLP on both macro ROC-AUC (0.9861 vs.\ 0.9859)
and macro PR-AUC (0.9520 vs.\ 0.9514)---an inversion to which an accuracy/F1 selection
rule is blind. Only the decision tree is clearly behind.

\begin{table}[t]
\centering
\caption{Model comparison on the holdout set (experimental result)}
\label{tab:modelcompare}
\begin{tabularx}{\linewidth}{@{}l C C C C@{}}
\toprule
Model & Accuracy & Macro-F1 & ROC-AUC & PR-AUC \\
\midrule
MLP (deployed) & 0.910 & 0.881 & 0.9859 & 0.9514 \\
Logistic regression & 0.892 & 0.864 & 0.9861 & 0.9520 \\
Random forest & 0.895 & 0.868 & 0.9764 & 0.9233 \\
Decision tree & 0.795 & 0.765 & 0.9276 & 0.8359 \\
\bottomrule
\end{tabularx}
\end{table}

\begin{figure}[t]
\centering
\includegraphics[width=0.72\linewidth]{fig_model_compare_v2.png}
\caption{Accuracy, macro-F1, ROC-AUC, and PR-AUC for the four classifiers. The MLP and
logistic regression are statistically indistinguishable on this dataset (see
Table~\ref{tab:stats}).}
\label{fig:modelcompare}
\end{figure}

\subsection{Statistical Validation}\label{subsec:statistical}
\textbf{\textit{Experimental result.}} Because every classifier was evaluated on the same
600-case holdout with an identical class ordering, model differences can be tested on
paired predictions. For each pairwise comparison the null hypothesis is that the two
models have equal population error rates---equivalently, that their discordant predictions
are symmetric---against the two-sided alternative of unequal error rates, at a significance
level of $\alpha = 0.05$. McNemar's test is the appropriate paired test for this design;
it is computed with a continuity correction, and for the small discordant counts here an
exact binomial test yields the same decisions. Table~\ref{tab:stats} reports each
comparison against the deployed MLP.

The MLP is not statistically distinguishable from logistic regression ($\chi^2 = 2.70$ on
24/13 discordant pairs, $p = 0.100$) or from the random forest ($\chi^2 = 1.83$ on 22/13
pairs, $p = 0.176$) at $\alpha = 0.05$; only the decision tree is beaten significantly
($\chi^2 = 48.67$ on 82/13 pairs, $p < 0.001$). The discordant-pair odds ratios---a
McNemar effect-size measure---are correspondingly modest for the two close baselines (1.85
and 1.69) and large only for the decision tree (6.31). Because three comparisons are made
against the same MLP, a Holm--Bonferroni correction is applied across them; it leaves
every decision unchanged, since the decision-tree result stays below its adjusted
threshold while the logistic-regression and random-forest results stay above theirs.

These tests are read alongside the interval estimates. The headline metrics carry a
bootstrap 95\% confidence interval computed over 300 resamples of the holdout: holdout
accuracy has a point estimate of 0.910 with a 95\% interval of [0.888, 0.933], and
macro-F1 a point estimate of 0.881 with a 95\% interval of [0.847, 0.912]. Five-fold
stratified cross-validation is consistent with these figures (accuracy
$0.912 \pm 0.009$, macro-F1 $0.867 \pm 0.018$). The overlapping intervals and the
non-significant tests together mean the evidence does not support a claim that a neural
model is necessary for this task; the MLP is retained for its probability outputs, which
the occlusion explainer consumes, and the non-significant comparisons are reported as an
absence of demonstrated superiority rather than as proof of equivalence.

\begin{table}[t]
\centering
\caption{Statistical validation of pairwise model differences (experimental result)}
\label{tab:stats}
\footnotesize
\setlength{\tabcolsep}{3pt}
\begin{tabularx}{\linewidth}{@{}L c c c c c L@{}}
\toprule
Comparison & Test & $\chi^2$ (b/c) & $p$-value & $\alpha$ & Significant? & Interpretation \\
\midrule
MLP vs.\ logistic regression & McNemar & 2.70 (24/13) & 0.100 & 0.05 & No & No demonstrated difference \\
MLP vs.\ random forest & McNemar & 1.83 (22/13) & 0.176 & 0.05 & No & No demonstrated difference \\
MLP vs.\ decision tree & McNemar & 48.67 (82/13) & $<$ 0.001 & 0.05 & Yes & MLP significantly better \\
\bottomrule
\end{tabularx}
\par\vspace{3pt}
{\footnotesize\itshape\raggedright $\chi^2$ is McNemar's statistic with continuity
correction and (b/c) are the discordant-pair counts---cases the MLP classifies correctly
and the comparator does not (b), versus the reverse (c). All $p$-values are two-sided; the
`Significant?' decision reflects the Holm--Bonferroni-adjusted threshold, which coincides
with the uncorrected decision here.\par}
\end{table}

\subsection{Guardrail Dominance and Reconciliation}\label{subsec:guardrail-dom}
\textbf{\textit{System output.}} Applied across the dataset, the guardrail alone assigns
5.87\% of users to Low, 33.53\% to Medium, and 60.60\% to High risk. The guardrail signal
is greater than or equal to the cluster signal for 100\% of users and strictly greater for
84.13\%, so the guardrail---not the classifier---is the binding signal in the
reconciliation for most users. Read the other way (Fig.~\ref{fig:guardrail}), when the
model predicts Low its prediction never changes the final band (0.00\% of cases), when it
predicts Medium the final band changes for 5.87\% of cases, and when it predicts High it
changes for 39.40\%. Substituting a constant Low-risk predictor would alter 0.00\% of
final decisions.

This is the single most consequential finding, and it qualifies a naive reading of the
91\% accuracy. Because the conservative guardrail escalates so many users to High risk on
cash-flow grounds alone, the learned classifier determines the end recommendation for only
a minority of users. The accuracy is a genuine property of the classifier, but it is not a
measure of the system's end-to-end decision quality, which is dominated by the
deterministic layer.

\begin{figure}[t]
\centering
\includegraphics[width=0.9\linewidth]{fig_guardrail_dominance.png}
\caption{Guardrail dominance. (a) Risk share under clustering alone versus the guardrail
alone; the guardrail escalates 60.6\% of users to High risk. (b) Share of cases in which
the model prediction can still change the final band, by predicted class.}
\label{fig:guardrail}
\end{figure}

\subsection{Portfolio Simulation}\label{subsec:portfolio-sim}
\textbf{\textit{Simulation result.}} Under the deterministic twelve-month projection with
fixed monthly returns, the corrected return on investment is 25.72\% for the Low-risk
allocation, 15.40\% for Medium, and 2.43\% for High, against a USDC-only baseline of
2.43\% (Table~\ref{tab:roi}, Fig.~\ref{fig:roi}). Maximum drawdowns are $-3.79\%$,
$-2.08\%$, and 0.00\% respectively. The High allocation and the baseline coincide exactly
because a High-risk user is placed entirely in USDC.

\begin{table}[t]
\centering
\caption{Original versus corrected twelve-month ROI by risk band}
\label{tab:roi}
\begin{tabularx}{\linewidth}{@{}l C C C C@{}}
\toprule
Risk band & BTC/ETH/USDC & Original ROI & Corrected ROI & Max drawdown \\
\midrule
Low & 35 / 35 / 30 & 20.18\% & 25.72\% & $-3.79\%$ \\
Medium & 20 / 20 / 60 & 12.34\% & 15.40\% & $-2.08\%$ \\
High & 0 / 0 / 100 & 2.22\% & 2.43\% & 0.00\% \\
\midrule
Baseline (USDC) & 0 / 0 / 100 & --- & 2.43\% & 0.00\% \\
\bottomrule
\end{tabularx}
\end{table}

\begin{figure}[t]
\centering
\includegraphics[width=0.72\linewidth]{fig_portfolio_roi.png}
\caption{Deterministic twelve-month simulated ROI by risk band: originally reported values
(inflated by the backtester indexing error) versus corrected values, against the USDC-only
baseline. The returns are fixed synthetic constants, so the figure demonstrates allocation
mechanics rather than market performance.}
\label{fig:roi}
\end{figure}

Two cautions attach. First, these are the corrected values: the originally reported ROIs
were inflated by an off-by-one normalization in the backtester. Second, and more
fundamentally, the returns are generated from fixed, hardcoded monthly constants rather
than live or historical market data, so they demonstrate the mechanics of the allocation
policy---higher assessed risk yields a more conservative, lower-variance allocation---but
are not evidence of real-world investment performance. They are best read as an internal
consistency check on the allocation logic.

\subsection{Auditability}\label{subsec:audit}
\textbf{\textit{Theoretical property.}} Every recommendation is hashed with SHA-256 over a
canonicalized payload and stored with a unique transaction identifier, and a verification
routine re-hashes a stored record to detect alteration. Tamper-evidence here is a property
of the cryptographic construction rather than a measured outcome: any change to a logged
payload changes its digest and is detectable on verification. The mechanism is exercised
against an in-memory simulated ledger, not a deployed blockchain, so it demonstrates the
auditability pattern rather than production-grade distributed immutability.
